\documentclass{article}
\usepackage[margin=1in]{geometry}
\usepackage{mathtools, amsfonts, amsthm, xcolor, listings}

\title{HW 05}
\author{Sam Ly}

\begin{document}
\maketitle

\section*{Exercise 1}
Collaborators: Sam

\noindent
Tools and resources used: I used a lot of stack exchange and a little Gemini.
\pagebreak

\section*{Exercise 2}
\begin{itemize}
  \item Points attempted: 3
  \item Self-assessed score: 3.
  \item Time spent: 3 minutes.
  \item Difficulty: 1/10.
  \item Questions/Feedback: None.
\end{itemize}
Done: For the blog post, I am looking to do something related to the topology 
prompt. I am currently only in a group of 2, and that is with Justin.
\pagebreak

\section*{Exercise 3}
\begin{itemize}
  \item Points attempted: 10
  \item Self-assessed score: 10.
  \item Time spent: 60 minutes.
  \item Difficulty: 6/10.
  \item Questions/Feedback: Correctness, clarity.
\end{itemize}

\begin{enumerate}
    \item[Problem 7.] {
        \begin{enumerate}
            \item[7.1] {
                If \(G\) is a group of infinite order, then all subgroups of \(G\)
                have infinite order.

                False: A counterexample would be that the group of \(\{1, -1\) under 
                multiplication is a subgroup of the group of non-zero rationals 
                under multiplication.
            }

            \item[7.2] {
                If \(G\) is an abelian group, then all subgroups of \(G\) are abelian 
                groups.

                True: Commutativity applies to all elements of \(G\), so by definition 
                a subgroup of the elements of \(G\) must also be commutative.
            }

            \item [7.3] {
                If \(H\) is a normal subgroup of \((G, *)\), then \(g * h = h * g\) 
                for all  \(g \in G\) and \(h \in H\).

                False: The normal subgroup invariant only requires that the cosets 
                \(gH = Hg\), however, it doesn't necessarily mean that \(g * h = h * g\), 
                just that \(g * h\) and \(h * g\) are in the same coset \(gH = Hg\).
            }

            \item [7.4] {
                Suppose that \((X, \star)\) is a set with a binary operation that 
                is not necesssarily a group, and there exists a ``left identity'' 
                \(x \in X\) such that \(x \star a = a\) for all \(a \in X\). Then 
                \(x\) is also a ``right identity'', that is \(a \star x = a\) 
                for all \(a \in X\).

                False: Imagine that the binary operation \(\star\) just selects the 
                element on its right, so all elements are the ``left identity'' 
                and no elements are the ``right identity.''
            }

            \item [7.5] {
                If \(G\) is a group and \(g \in G\), then every column of the 
                Cayley table for \(G\) contains \(g\) exactly once.

                True: we know this to be true because for every \(g' \in G\), 
                \(g * g^{-1} * g' = g'\), so that must be a column of \(g\). 
            }

            \item [7.6] {
                The set of rational numbers \(\mathbb{Q}\) forms a group under 
                multiplication.

                False: There is no inverse of 0.
            }

            \item [7.7] {
                If \(G\) is a finite group of order \(n\), then for every element \(g \in G\),
                \(g^n = e\).

                True: the order \(m\) of all subgroups \(H\) must be a factor of \(G\), 
                so considering all subgroups generated by a single element \(g \in G\), 
                the order of \(\langle g \rangle\) must be a factor of the order of 
                \(G\), so \(g^n = e^{n/m} = e\).
            }
        \end{enumerate}

        \item [Problem 8] {
            Let \(G\) be a (not necessarily abelian) group, and let \(H\) be a 
            (not necessarily) subgroup of \(G\). Let the normalizer of \(H\) be 
            defined as 

            \[N_G(H) = \{ g \in G \mid gH = Hg \} .\]

            \begin{enumerate}
                \item [8.1] {
                    Prove that \(N_G(H) \le G\).

                    \begin{proof}
                        We begin by noticing that \(N_G(H)\) is a subset of 
                        \(G\). Thus, the group axiom of associativity, is 
                        already satisfied. 
                        
                        First we show that \(N_G(H)\) has an identity element.
                        Notice \(e_G \in N_G(H)\) because \(e_G H = H e_G\). So, 
                        \(e_G\) is the identity element of \(N_G(H)\).

                        Next, assuming that \(n \in N_G(H)\), we can show that 
                        \(n^{-1}\) must also be in \(N_G(H)\).  First, we see that 
                        for every \(n\) and \(h_1 \in H\), there exists a unique 
                        \(h_2 \in H\), such that \(n * h_1 = h_2 * n\). Now, we 
                        left and right multiply by \(n^{-1}\), to get 
                        \begin{align}
                            n * h_1 &= h_2 * n \\
                            n^{-1} * n * h * n^{-1} &= n^{-1} * h_2 * n * n^{-1} \\
                            h * n^{-1} &= n^{-1} * h_2.
                        \end{align}
                        Therefore, \(n^{-1} \in N_G(H)\), so every element has an 
                        inverse.

                        Lastly, we show that \(N_G(H)\) is closed. That is, 
                        for all \(n_1, n_2 \in N_G(H)\), \(n_1 * n_2 \in N_G(H)\) 
                        and \((n_1 * n_2)H = H(n_1 * n_2)\).
                        Let us pick some \(h_1, h_1', h_2, h_2' \in H\), such that 
                        \begin{align}
                            n_1 * h_1 &= h_1' * n_1 \\
                            n_2 * h_2 &= h_2' * n_2.
                        \end{align}

                        We isolate \(n_1\)  in equation (4) as \(n_1 = h_1' * n _1 * h_1^{-1}\).
                        We left multiply equation (5) by \(n_1\) to get
                        \begin{align}
                            n_1 * n_2 * h_2 &= h_1' * n _1 * h_1^{-1} * h_2' * n_2.
                        \end{align}
                        Now, because \(n_1 H = H n_1\), there must exist some 
                        \(h_1'^{-1}\) where \(n_1 * h_1^{-1} = h_1'^{-1} * n_1\).
                        We substitute this into equation (6) to get 
                        \begin{align}
                            &= h_1' * h_1'^{-1} * n_1 * h_2' * n_2.
                        \end{align}
                        Similarly for \(n_1 * h_2'\), there exists some \(h_2''\) 
                        where \(n_1 * h_2' = h_2'' * n_1\). So, we have,
                        \begin{align}
                            &= h_1' * h_1'^{-1} * h_2'' * n_1 * n_2.
                        \end{align}

                        Since \(H\) itself is closed, \(h_1' * h_1'^{-1} * h_2'' \in H\).
                        Therefore, \(n_1 * n_2\) is also a normalizer \(N_G(H)\).
                        Thus, \(N_G(H)\) is is closed, so \(N_G(H)\) is a subgroup 
                        of \(G\).
                    \end{proof}
                }

                \item [8.2] {
                    Prove that for all \(g \in N_G(H)\) and \(h \in H\), 
                    \(g h g^{-1} \in H\).

                    \begin{proof}
                        We begin by noticing that \(gH = Hg\). That is, there 
                        for every \(h \in H\), there exists a unique \(h' \in H\)
                        where \(g * h = h' * g\). We right multiply by \(g^{-1}\) 
                        to get \(g * h * g^{-1} = h'\). Thus, for all 
                        \(g \in N_G(H)\) and \(h \in H\), \(g h g^{-1} \in H\). 
                    \end{proof}
                }
            \end{enumerate}
        }
    }
\end{enumerate}


\end{document}
